\subsection{R Symmetry Breaking Versus Supersymmetry Breaking --- arXiv:hep-ph/9309299}

\textbf{Authors:} Ann E. Nelson, Nathan Seiberg

\paragraph{主な主張}
一般的な超ポテンシャルをもつ計算可能なWess--Zumino有効理論では、SUSY破れにR対称性が必要であり、その自発的破れはSUSY破れの十分条件になることを示した\cite{Nelson:1993nf}。

\paragraph{新規性}
R対称性とSUSY破れの必要・十分条件を明確化し、一般性の仮定を外すとR対称性を明示的に破ってもSUSY破れを保てる可能性を示した。

\paragraph{位置付け}
SUSY破れ模型とR-axionの問題を整理するNelson--Seibergの議論の原論文。

\paragraph{コメント（読書メモ）}
R対称性とSUSY breakingの関係を示す原論文として参照。
